% -*- coding: utf-8 -*-
% !TEX program = xelatex
\documentclass[plain,noanswer,medmath]{../../template/jnuexam}
\usepackage{../../template/kysx}

\begin{document}


\examtitle{1998}{4}


\loadproblems{3}{1998P3}%载入试卷三


\makepart{填空题}{本题共5小题，每小题3分，满分15分}


\useproblem{3}{1}{1}%【同试卷三第一(1)题】


\useproblem{3}{1}{2}%【同试卷三第一(2)题】


\useproblem{3}{1}{4}%【同试卷三第一(4)题】


\begin{problem}
设$A$, $B$均为$n$阶矩阵，$|A| = 2$, $|B| = - 3$，则$\left| 2A^*B^{-1} \right| =$ \fillin{}．
\end{problem}


\begin{problem}
设一次试验成功的概率为$p$，进行$100$次独立重复试验，当$p=$\fillin{}时，
成功次数的标准差的值最大；其最大值为 \fillin{}．
\end{problem}


\makepart{选择题}{本题共5小题，每小题3分，满分15分}


\useproblem{3}{2}{1}%【同试卷三第二(1)题】


\useproblem{3}{2}{2}%【同试卷三第二(2)题】


\begin{problem}
若向量组$\alpha ,\beta ,\gamma$线性无关；$\alpha ,\beta ,\delta$线性相关，则\pickout{}
\begin{abcd}
\item $\alpha$必可由$\beta ,\gamma ,\delta$线性表示．
\item $\beta$必不可由$\alpha ,\gamma ,\delta$线性表示．
\item $\delta$必可由$\alpha ,\beta ,\gamma$线性表示．
\item $\delta$必不可由$\alpha ,\beta ,\gamma$线性表示．
\end{abcd}
\end{problem}


\begin{problem}
设$A$, $B$, $C$是三个相互独立的随机事件，且$0 < P(C) < 1$，则在下列给定的四对事件中不相互独立的是\pickout{}
\begin{abcd}
\item $\overline{A + B}$与$C$．
\item $\overline{AC}$与$\bar C$．
\item $\overline{A - B}$与$\bar C$．
\item $\overline{AB}$与$\bar C$．
\end{abcd}
\end{problem}


\useproblem{3}{2}{5}%【同试卷三第二(5)题】


\makepart{}{本题满分6分}

\begin{problem}
求$\lim_{n \to \infty} \left(n\tan \frac{1}{n} \right)^{n^2}$（$n$为自然数）．
\end{problem}


\makepart{}{本题满分6分}%【同试卷三第三题】

\useproblem{3}{3}{0}


\makepart{}{本题满分5分}%【同试卷三第四题】

\useproblem{3}{4}{0}


\makepart{}{本题满分6分}%【同试卷三第五题】

\useproblem{3}{5}{0}


\makepart{}{本题满分6分}

\begin{problem}
设函数$f(x)$在$\left[a,b \right]$上连续，在$\left(a,b \right)$内可导，且$f(a) = f(b) = 1$，
试证存在$\xi ,\;\eta \in(a,b)$，使得$\e^{\eta - \xi}[f(\eta) + f'(\eta)] = 1$．
\end{problem}


\makepart{}{本题满分9分}

\begin{problem}
设直线$y = ax$与抛物线$y = x^2$所围成图形的面积为$S_1$，
它们与直线$x = 1$所围成的图形面积为$S_2$，并且$a < 1$．
\begin{enumerate}
\item 试确定$a$的值，使$S_1 + S_2$达到最小，并求出最小值；
\item 求该最小值所对应的平面图形绕$x$轴旋转一周所得旋转体的体积．
\end{enumerate}
\end{problem}


\makepart{}{本题满分9分}%【同试卷三第九题】

\useproblem{3}{9}{0}


\makepart{}{本题满分9分}

\begin{problem}
已知下列非齐次线性方程组①和②：
$$\text{①}\;\begin{cases}
  x_1 + x_2 - 2x_4 = - 6, \\
  4x_1 - x_2 - x_3 - x_4 = 1, \\
  3x_1 - x_2 - x_3 = 3;
\end{cases}\qquad\text{②}\;\begin{cases}
  x_1 + mx_2 - x_3 - x_4 = - 5, \\
  nx_2 - x_3 - 2x_4 = - 11, \\
  x_3 - 2x_4 = - t + 1.
\end{cases}$$
\begin{enumerate}
\item 求解方程组①，用其导出组的基础解系表示通解．
\item 当方程组②中的参数$m$, $n$, $t$为何值时，方程组①与②同解？
\end{enumerate}
\end{problem}


\makepart{}{本题满分7分}

\begin{problem}
求某种商品每周的需求量 X 是服从区间$\left[10,30 \right]$上均匀分布的随机变量，
而经销商进货数量为区间$\left[10,30 \right]$中的某一整数，商店每销售一单位商品可获利$500$元；
若供大于求则削价处理，每处理$1$单位商品亏损$100$元；若供不应求，则可从外部调剂供应，
此时每$1$单位商品仅获利$300$元，为使商品所获利润期望值不小于$9280$元，试确定最少进货量．
\end{problem}


\makepart{}{本题满分7分}

\begin{problem}
某箱装有$100$件产品，其中一、二、三等品分别为$80$件、$10$件和$10$件，
现在从中随机抽取一件，记$X_i = \begin{cases}
  1, & \text{若抽到$i$等品} \\
  0, & \text{其他}
\end{cases}$($i = 1,2,3$)，试求：
\begin{enumerate}
\item 随机变量$X_1$与$X_2$的联合分布；
\item 随机变量$X_1$与$X_2$的相关系数$\rho$．
\end{enumerate}
\end{problem}


\end{document}
